\chapter{시간을 지키는 법: 시점 학습과 편향}

2장에서 수익률의 자를 정했다. 이 장은 그 자로 \textbf{언제의 정보로} 학습하고 \textbf{어떤 대출로} 채점하는지를 다룬다. 누수는 컬럼(1장)에서만 생기지 않는다. 행에서도 생긴다.

\section{무작위 분할은 왜 성적을 부풀리나}
보통 머신러닝 수업에서는 데이터를 무작위로 80:20으로 나눈다. LC 데이터에 그렇게 하면 2016년 대출로 학습해 2013년 대출을 맞히는 일이 생긴다. 2016년의 경기, LC의 심사 기준, 금리 수준은 2013년 투자자가 알 수 없던 정보다.

\begin{keybox}[결과 C5]
무작위 분할 AUC 약 0.73, 시점 학습 AUC 0.62--0.68. 차이 전부가 ``미래를 본 값''이다.
\Intuit{내년 기출문제로 공부하고 올해 시험을 보면 점수는 높지만, 실제로는 그런 일이 불가능하다.}
\end{keybox}

\section{시점 학습(as-of training)}
연도 $Y$에 발행된 대출을 평가할 때, 모델은 \textbf{$Y$년 1월 1일에 이미 결과가 확정된 대출}로만 학습한다(\ft{asof.training_set}). 확정 조건은 두 가지다.

\begin{enumerate}[leftmargin=*]
\item 그 날짜 전에 발행되었다.
\item 그 날짜까지 결과가 났다. 정상 상환은 마지막 납부 달에, 부도는 \textbf{마지막 납부 + 5개월}에 확정된 것으로 본다(\ft{CHARGE_OFF_LAG = 5}).
\end{enumerate}

\begin{figure}[H]\centering
\begin{tikzpicture}[x=1.0cm,y=0.75cm,font=\small]
\draw[-{Stealth}] (0,0)--(9.6,0) node[right]{시간};
\foreach \x/\l in {1/2010,3/2011,5/2012,7/2013,9/2014} {\draw (\x,0.1)--(\x,-0.1) node[below]{\l};}
\draw[very thick] (7,-0.6)--(7,4.3) node[above]{기준일 2013-01-01};
\draw[black!70,line width=2pt] (1.2,1)--(4.5,1); \node[right] at (9.8,1) {A: 정상 상환, 확정};
\draw[black!70,line width=2pt] (2,2)--(5.4,2); \draw[black!35,line width=2pt,dashed] (5.4,2)--(6.4,2); \node[right] at (9.8,2) {B: 부도, 5개월 뒤 확정};
\draw[black!70,line width=2pt] (4,3)--(6.6,3); \draw[black!35,line width=2pt,dashed] (6.6,3)--(7.6,3); \node[right] at (9.8,3) {C: 기준일 뒤 확정, 제외};
\draw[black!70,line width=2pt] (5.5,4)--(9.5,4); \node[right] at (9.8,4) {D: 진행 중, 제외};
\end{tikzpicture}
\caption{시점 학습의 학습 표본. 2013년 대출을 평가할 모델은 A와 B로만 학습한다. 점선은 상각 지연 5개월.}
\end{figure}

\begin{whybox}[왜 5개월인가]
LC는 연체 120일(약 4개월)이 넘으면 상각(charge-off)한다. 마지막 납부 이후 상각이 장부에 잡히기까지 약 5개월이 걸린다. 지연을 0으로 두면 아직 ``연체 중''이던 대출을 이미 부도로 알았던 것처럼 학습해 미래 정보가 섞인다. 6단계에서 지연 0과 5를 비교했고, 가정에 따라 편향 방향이 바뀌는 것을 확인했다.
\end{whybox}

\paragraph{대가.} 시점 학습은 표본이 작다. 2013년 모델의 학습 대출은 11,345건이고, 2016년 모델은 156,153건이다. 이른 해일수록 모델이 약하고 불안정하다(5장의 2013년 결과들이 흔들리는 이유).

\section{종결 표본 편향(resolution bias)}
\textbf{질문.} 평가할 때 ``이미 결과가 난 대출''만 쓰면 무엇이 달라지나?

\begin{workout}[4건짜리 종결 편향]
\Use{censoring, 종결 표본 편향(이 절)}
\Step{1} 같은 달 발행된 36개월 대출 4건(예시):
A 저금리 6\%, 3년 내내 정상. B 고금리 20\%, 3년 내내 정상. C 고금리 20\%, 6개월 뒤 부도. D 저금리 6\%, 1년 안에 조기상환.
\Step{2} 3년 뒤 전체를 보면 부도율 $1/4=25\%$. 고금리 B의 이자가 C의 손실을 메운다.
\Step{3} 발행 1년 뒤에 ``끝난 대출만'' 보면 C와 D만 남는다. 부도율 $1/2=50\%$. \Concept{censoring}
\Step{4} 남은 표본에서 ``고금리를 고르는 규칙''은 C만 고르게 되어 최악이고, ``부도확률이 낮은 것을 고르는 규칙''은 D를 골라 최선이다.
\Ans{일찍 본 표본은 부도를 과대 대표하고, 오래 이자를 받아야 이득인 고금리 정상 대출(B)을 숨긴다.}
\Trap{``종결 대출만 쓰면 깨끗하다''고 생각한다. 깨끗한 것은 라벨이고, 표본 구성은 오히려 왜곡된다.}
\end{workout}

\subsection*{실제 데이터로 시뮬레이션(26단계)}
2013--15년 36개월 대출은 지금 모두 끝나 정답을 안다. 같은 대출을 관측 시점에 맞춰 일부러 잘라, 두 순위 방식의 선별 가치(6장)를 비교했다. LightGBM 시드 10개, 발행월 부트스트랩 1만 번, 판정 기준은 실행 전에 고정했다.

\begin{table}[H]\centering\small
\begin{tabular}{@{}lrrrrr@{}}\toprule
관측 & 남은 비율 & 부도율 & 기대수익 순위 & 부도확률 순위 & 차이(95\% 구간)\\\midrule
발행 후 1--2년 & 약 30\% & 약 18\% & 1.44\%p & 5.28\%p & $-3.83$ ($-4.67$, $-3.04$)\\
발행 후 2--3년 & 약 59\% & 약 18\% & 1.18\%p & 2.21\%p & $-1.03$ ($-1.36$, $-0.70$)\\
전체 & 100\% & 약 14\% & 1.01\%p & 0.10\%p & $+0.91$ ($+0.73$, $+1.10$)\\\bottomrule
\end{tabular}
\caption{종결 편향 시뮬레이션(\ft{scripts/26_resolution_bias.py}). 단위는 무작위 승인 대비 \%p.}
\end{table}

\begin{keybox}[결과 C23과 C6 정정]
관측 기간이 짧을수록 부도확률 순위가 이기는 것처럼 보이고, 전체를 보면 기대수익 순위가 이긴다. 14단계의 2017년 보류 구간에서 부도확률 순위가 이긴 것은 이 편향으로 설명된다.
\Intuit{일찍 채점하면 ``빨리 망한 대출''만 보이므로, 망하지 않는 것만 신경 쓰는 규칙이 똑똑해 보인다.}
\end{keybox}

\begin{warnbox}[고칠 수 있나]
종결 여부는 결과(부도, 조기상환) 자체에 달려 있다. 이런 결측은 MNAR(4장)이라, 특성으로 만든 가중치(inverse probability weighting)로는 고칠 수 없다. 대신 \textbf{test를 보기 전에 읽는 법을 정했다}: 발행연도별 만기 도달 비율부터 확인하고, 비율이 높은 연도만 두 순위 비교의 근거로 쓴다.
\end{warnbox}

\section{생존분석은 왜 안 썼나}
진행 중 대출 62만 건을 버리는 것이 아깝다면, 생존분석(survival analysis)으로 ``아직 안 끝난'' 정보까지 쓸 수 있다. 시도했지만 단순 PD 모델을 이기지 못했다(연구 노트 실패 목록). 수준(확률 크기)에는 편향이 있어도 순위는 크게 무너지지 않았기 때문이다.

\begin{career}
시계열이 섞인 데이터에서 무작위 분할을 쓰지 않은 이유, 그리고 \textbf{평가 표본의 선택 편향}까지 시뮬레이션으로 정량화한 것은 강한 이야깃거리다. ``모델 성능을 올렸다''보다 ``성능이 부풀려지는 경로 두 가지(무작위 분할, 종결 편향)를 찾아 크기를 쟀다''가 훨씬 드물다.
\end{career}
